%%%%%%%%%%%%%%%%%%%%%%%%%%%%%%%%%%%%%%%%%
% This template has been downloaded from:
% http://www.LaTeXTemplates.com
%
% Original author:
% Rensselaer Polytechnic Institute (http://www.rpi.edu/dept/arc/training/latex/resumes/)
%
% Important note:
% This template requires the res.cls file to be in the same directory as the
% .tex file. The res.cls file provides the resume style used for structuring the
% document.
%
%%%%%%%%%%%%%%%%%%%%%%%%%%%%%%%%%%%%%%%%%

%----------------------------------------------------------------------------------------
%	PACKAGES AND OTHER DOCUMENT CONFIGURATIONS
%----------------------------------------------------------------------------------------
\documentclass[margin, 10pt]{res} % Use the res.cls style, the font size can be changed to 11pt or 12pt here

%\usepackage{helvet} % Uncomment to use the Helvetica postscript font; the CV uses the default Computer Modern font
%\usepackage{newcent} % To change the default font to the new century schoolbook postscript font uncomment this line and comment the one above
\usepackage{hyperref}
%\usepackage{url}
\usepackage{enumitem}

%\setlength{\textwidth}{5.1in} % Text width of the document
\setlength{\textwidth}{5.3in} % Text width of the document

%\resumewidth=10in
%\sectionwidth=1in
\sectionskip=0.3in
\begin{document}

\moveleft.5\hoffset\centerline{\Large\bf Zhiyuan Yao} % Your name at the top
\smallskip
\moveleft.5\hoffset\centerline{\large Postdoctoral Researcher, Hebrew University of Jerusalem}
\smallskip
\moveleft.5\hoffset\centerline{\large Curriculum Vitae}

%\moveleft\hoffset\vbox{\hrule width\resumewidth height 1pt}\smallskip % Horizontal line after name; adjust line thickness by changing the '1pt'
%----------------------------------------------------------------------------------------
\begin{resume}
%----------------------------------------------------------------------------------------
%     CONTACT INFORMATION
%----------------------------------------------------------------------------------------

\section{Contact \\ Information}

\begin{minipage}{0.55\textwidth}
  Office 209, Ross building\\
  Hebrew University of Jerusalem \\
  Givat Ram, 9190401 \\
  Jerusalem, Israel
\end{minipage}
\begin{minipage}{0.55\textwidth}
  \begin{tabular}{r}
    +972-58-475-3627 \\
   zhiyuan.yao@mail.huji.ac.il \\
    ORCID: 0000-0003-1604-1272 \\
  \end{tabular}
\end{minipage}

%----------------------------------------------------------------------------------------
%	EDUCATION SECTION
%----------------------------------------------------------------------------------------
\section{Education}

\textbf{Postdoc, Hebrew University of Jerusalem}, Israel\\
10/2022 - Present\\
Advisor: Nir Mandelker, Avishai Dekel\\
\\
\textbf{Joint Ph.D, Max Planck Institute for Astrophysics}, Germany\\
05/2021 - 03/2022\\
Advisor: Volker Springel\\
\\
\textbf{Ph.D, Astronomy, Shanghai Astronomical Observatory}, Chinese Academy of Sciences, China\\
09/2016 - 06/2022 \\
Advisor: Feng Yuan\\
\\
\textbf{Bachelor, Astronomy, University of Science and Technology of China}, China\\
09/2012 - 06/2016

%----------------------------------------------------------------------------------------
%	RESEARCH INTERESTS
%----------------------------------------------------------------------------------------
\section{Research \\ Interests}

Galaxy formation and evolution; circumgalactic medium; multiphase hydrodynamics; cold streams and virial shocks; turbulence and mixing; numerical simulations; high-redshift galaxies; computational astrophysics

%----------------------------------------------------------------------------------------
%	Honors & Awards
%----------------------------------------------------------------------------------------
\section{Honors \& \\ Awards}

2024, Rosenblum prize\\
2021-2022, Sino-German Joint PhD Exchange Program (Shanghai Astronomical Observatory and Max Planck Institute for Astrophysics, 04/2021 - 03/2022)\\
2020, Academic Scholarship of University of Chinese Academy of Sciences

%----------------------------------------------------------------------------------------
%	Publications
%----------------------------------------------------------------------------------------
\section{Selected \\ Publications}

\textbf{First-author publications}\\[4pt]
\textbf{Yao Z.}, Mandelker N., Oh S. P.\\
``Cold stream penetration of virial shocks'', submitted to MNRAS, arXiv:2607.14090 (2026)\\[6pt]
\textbf{Yao Z.}, Mandelker N., Oh S. P., Aung H., Dekel A.\\
``Effects of cloud geometry and metallicity on shattering and coagulation of cold gas, and implications for cold streams penetrating virial shocks'', MNRAS, 536, 3053 (2025)\\[6pt]
\textbf{Yao Z.}, Yuan F., Ostriker J. P.\\
``Active galactic nucleus feedback in an elliptical galaxy with updated AGN physics: parameter exploration'', MNRAS, 501, 398 (2021)\\[6pt]
\textbf{Yao Z.} \& Gan Z.\\
``Hot gas flows on a parsec scale in the low-luminosity active galactic nucleus NGC 3115'', MNRAS, 492, 444 (2020)\\[10pt]
\textbf{Collaborative publications}\\[4pt]
Dekel, A. et al.\\
``From Feedback-Free Star Clusters to Little Red Dots via Compaction'', arXiv:2511.07578 (2025)\\[6pt]
Dekel, A. et al.\\
``From FFB Starbursts at Cosmic Dawn to Quenching at Cosmic Morning: Hi-z Galaxy Bimodality'', MNRAS, 544, 160 (2025)\\[6pt]
Yang, G. et al.\\
``Do current X-ray observations capture most of the black-hole accretion at high redshifts?'', ApJ, 921, 170 (2021)

%----------------------------------------------------------------------------------------
%	Invited Talks
%----------------------------------------------------------------------------------------
\section{Invited Talks}

12/2025, Tel Aviv University Seminar, Tel Aviv, Israel\\
10/2021, MPA Institute Seminar, Garching, Germany

%----------------------------------------------------------------------------------------
%	Conference Contributions
%----------------------------------------------------------------------------------------
\section{Conference \\ Contributions}

11/2025, Modelling of Multiphase Astrophysical Media, Ringberg, Germany\\
08/2025, Galaxy Formation Workshop, Santa Cruz, United States\\
07/2025, Cosmic Ecosystems, Waterloo, Canada\\
08/2024, Galaxy Formation Workshop, Santa Cruz, United States\\
05/2023, Modelling of Multiphase Astrophysical Media, Kochel, Germany\\
08/2022, Galaxy Formation Workshop, Santa Cruz, United States\\
07/2021, Ringberg Retreat, Munich, Germany\\
Earlier (2018--2020): 4 contributed talks at meetings in China and Germany.

%----------------------------------------------------------------------------------------
%	Skills
%----------------------------------------------------------------------------------------
\section{Technical \\ Skills}

Hydrodynamic simulation codes: AREPO, RAMSES, ZEUS--MP/2\\
Programming: Python, C, C++, Fortran, MPI/OpenMP\\
Analysis tools: yt, NumPy, SciPy, Matplotlib\\
High-performance computing: distributed-memory simulations on large HPC systems

%----------------------------------------------------------------------------------------
%	Mentoring
%----------------------------------------------------------------------------------------
\section{Mentoring}

2025 -- Aviad Leibovich (B.Sc. student), Hebrew University of Jerusalem

%----------------------------------------------------------------------------------------
%	Professional Service
%----------------------------------------------------------------------------------------
\section{Professional \\ Service}

Organizer, Astrolunch Seminar Series, Hebrew University of Jerusalem

%----------------------------------------------------------------------------------------
%	References
%----------------------------------------------------------------------------------------
\section{References}

Nir Mandelker --- Hebrew University of Jerusalem --- nir.mandelker@mail.huji.ac.il\\[4pt]
Feng Yuan --- Fudan University --- fyuan@fudan.edu.cn\\[4pt]
Volker Springel --- Max Planck Institute for Astrophysics --- vspringel@mpa-garching.mpg.de

\end{resume}
\end{document}
